\documentclass[11pt]{article}
\usepackage[margin=1in]{geometry}
\usepackage{booktabs}
\usepackage{graphicx}
\usepackage{hyperref}
\usepackage{microtype}
\usepackage{amsmath}
\usepackage{xcolor}

\newcommand{\system}{\textsc{K12SimWorld}}
\newcommand{\benchmark}{\textsc{K12SimBench}}
\newcommand{\tbd}[1]{\textcolor{red}{[TBD: #1]}}

\title{K12SimWorld: Generating Executable and Pedagogically Grounded\\
Simulations for K--12 Multimodal Problems}
\author{Anonymous Authors}
\date{}

\begin{document}
\maketitle

\begin{abstract}
Multimodal models are increasingly evaluated on answers to educational questions, but an
answer does not reveal whether a model has formed a testable account of the depicted process.
We introduce \benchmark, a benchmark of approximately 240 Chinese K--12 multimodal problems
centered on physics with controlled cross-subject extensions, and \system, a framework that
turns a problem into an interleaved explanation and deterministic executable simulation.
The method anchors every scene to one typed world specification, making objects, parameters,
constraints, and expected events inspectable. We compare text-only reasoning, static diagrams,
direct code generation, unanchored simulation, and our full method using answer, event,
constraint, pedagogical, reliability, and cost measures. \tbd{Insert measured sample counts,
confidence intervals, and the single central empirical finding; do not claim learning gains.}
\end{abstract}

\section{Introduction}
K12Vista demonstrates the breadth and difficulty of Chinese K--12 multimodal reasoning
\cite{li2025k12vista}. EduIllustrate shifts evaluation from answers to interleaved text--diagram
explanations \cite{bi2026eduillustrate}, while VisPhyWorld treats executable simulator code as a
falsifiable physical hypothesis \cite{liang2026visphyworld}. Their combination is non-trivial:
K12Vista provides questions rather than reference videos, and many educational processes are
not rigid-body reconstruction tasks.

We ask whether models can generate an executable process that is simultaneously correct,
consistent across explanatory steps, and appropriate for a learner's grade. Our contributions
are: (1) a curated benchmark and expert-reference subset; (2) a canonical EduWorldSpec and
state-anchored generation protocol; and (3) a success-aware evaluation separating solution,
simulation, and pedagogical quality.

\section{Related Work}
We organize prior work around multimodal K--12 assessment, multimedia educational content,
programmatic diagram generation, and executable physical reasoning. Unlike pixel-space video
generation, our goal is not photorealism; unlike static diagram generation, our output commits
to state transitions and executable constraints. Multimedia-learning theory motivates evaluating
coordination and cognitive load rather than visual appeal alone \cite{mayer2002multimedia,
sweller1988cognitive,paivio1971imagery}.

\section{K12SimBench}
\paragraph{Source and selection.}
We curate problems from K12Vista without exposing reference solutions to candidate models.
The main subset targets middle- and high-school physics; mathematics, chemistry, biology, and
geography form a controlled extension. Inclusion requires pedagogically useful dynamics, a
clear answer and stepwise reference, a representable state, and deterministic short-horizon
execution. Knowledge-point groups remain disjoint across development, validation, and test.

\paragraph{Expert references.}
A stratified 60-problem subset receives canonical object states, parameter ranges, expected
events, final states, and verified reference simulations. \tbd{Report adjudication process,
license review, and final distributions from the frozen curation summary.}

\section{K12SimWorld}
Given problem $q$ and image $I$, the model produces a storyboard $B$, canonical world
specification $W$, and scene programs $C_1,\ldots,C_n$:
\[
(q,I) \rightarrow B \rightarrow W \rightarrow \{C_i(W)\}_{i=1}^{n}
\rightarrow \{\hat X_i\}_{i=1}^{n}.
\]
The storyboard alternates explanatory text and simulations only when motion or state change
supports a reasoning step. EduWorldSpec declares coordinates, objects, parameters with units,
constraints, initial and final state, expected events, learning goals, and visual conventions.
Every program records the canonical SHA-256 hash of $W$ and uses identical object identities,
values, and camera conventions. Mechanics routes to Three.js+Cannon.js, deterministic state
processes to P5.js, and symbolic construction to Manim. Validation permits one repair; a second
failure remains a scored failure rather than falling back to an empty scene.

\section{Evaluation}
We retain eight adapted EduIllustrate dimensions: correctness, coherence, pedagogy, typography,
simulation--problem alignment, layout, temporal consistency, and text--simulation coordination.
Four executable dimensions measure initial state, key events, final state, and constraint
satisfaction. We report separate geometric means $Q_{sol}$, $Q_{sim}$, and $Q_{ped}$, followed by
$Q_{all}=(Q_{sol}Q_{sim}Q_{ped})^{1/3}$. Failed generations receive zero on every dimension.

Each expert artifact is independently scored by at least three qualified raters. We report
Krippendorff's $\alpha$, per-dimension Spearman correlation with the automatic judge, paired
bootstrap confidence intervals, paired randomization tests, effect sizes, and Holm correction.

\section{Experimental Setup}
We compare text-only chain of thought, EduIllustrate-style static Manim, direct executable code,
independently generated scenes, and the full state-anchored system. Each frozen model--method
combination runs once on the full benchmark and three times on the expert subset. Prompt,
model/API version, renderer, repair budget, runtime, tokens, and contemporaneous price are logged.
\tbd{Insert exact model snapshots and inference dates.}

\section{Results}
\IfFileExists{../results/main_results.tex}{\input{../results/main_results.tex}}{
\begin{center}\fbox{Run the evaluation reporter to generate main\_results.tex.}\end{center}}

We answer seven preregistered questions: whether execution improves reasoning; whether gains are
pedagogical or merely visual; whether runnable code masks wrong parameters; the contributions of
WorldSpec and anchoring; the role of physics-native engines; difficult strata; and judge validity.
All claims in this section must name the supporting metric and confidence interval.

\section{Ablations and Analysis}
We remove WorldSpec, remove state anchoring, replace the physics-native backend, remove structured
visual context, and vary the repair budget. Results are paired by problem. Case studies include
one fully correct example, one executable but causally wrong example, one visually plausible but
constraint-violating example, and one pedagogically poor example.

\section{Limitations and Ethics}
Executable simulations necessarily simplify the world and may encode a plausible but incorrect
normalization when a problem omits parameters. Automated judges may be unreliable on subjective
pedagogy. The benchmark measures artifact quality, not student learning; without a classroom or
pre/post intervention we make no causal learning claim. Public release must respect K12Vista's
license and remove rater identity. Generated material requires qualified teacher review before use.

\section{Conclusion}
\system provides a falsifiable bridge between multimodal educational reasoning and programmatic
simulation. \tbd{Conclude only with findings supported by the completed experiment.}

\bibliographystyle{plain}
\bibliography{references}

\appendix
\section{Reproducibility Checklist}
Release the frozen IDs, schemas, prompt version, raw responses where permitted, validation logs,
renderer configuration, annotations, aggregation code, and exact API metadata. Do not release
provider credentials or identifiable rater data.

\section{Rubrics and Additional Cases}
Include complete 0--5 anchors, per-subject distributions, all ablation tables, failure taxonomy,
and representative cases selected before inspecting model identities.

\end{document}
